\section{Lecture 4: Learning as Optimization}

\subsection{Motivation}

In previous lectures, we developed three key ideas:
\begin{itemize}
    \item Learning as function approximation
    \item Learning as statistical estimation
    \item The role of hypothesis spaces and inductive bias
\end{itemize}

\medskip

\noindent
These ideas naturally lead to a central computational question:

\medskip

\noindent
\textit{How do we actually find the best function within a hypothesis space?}

\medskip

\noindent
The answer is to formulate learning as an optimization problem.

\medskip

\noindent
\textbf{Guiding idea.}  
Learning is the process of minimizing an objective function defined by data and model assumptions.

\subsection{Optimization Formulations of Learning}

Given a hypothesis space $\mathcal{H} = \{f_\theta\}$ and a loss function $\ell$, we define the empirical objective:
\[
L(\theta) = \frac{1}{n} \sum_{i=1}^n \ell(f_\theta(x_i), y_i).
\]

\medskip

\noindent
\textbf{Learning problem.}
\[
\min_{\theta \in \Theta} L(\theta).
\]

\medskip

\noindent
\textbf{Regularized formulation.}
\[
\min_\theta \; L(\theta) + \lambda \Omega(\theta),
\]
where $\Omega(\theta)$ is a regularization term.

\medskip

\noindent
\textbf{Examples.}
\begin{itemize}
    \item Linear regression: quadratic objective
    \item Logistic regression: convex but nonlinear objective
    \item Neural networks: highly non-convex objective
\end{itemize}

\medskip

\noindent
\textbf{Interpretation.}
\begin{itemize}
    \item Objective function encodes both data fit and inductive bias
    \item Learning reduces to searching for optimal parameters
\end{itemize}

\subsection{Gradient-Based Methods}

Many learning problems are solved using gradient-based optimization.

\medskip

\noindent
\textbf{Gradient descent.}
\[
\theta_{t+1} = \theta_t - \eta \nabla L(\theta_t),
\]
where $\eta > 0$ is the step size.

\medskip

\noindent
\textbf{Interpretation.}
\begin{itemize}
    \item The gradient points in the direction of steepest increase
    \item We move in the opposite direction to decrease the loss
\end{itemize}

\medskip

\noindent
\textbf{First-order approximation.}
\[
L(\theta + \delta) \approx L(\theta) + \nabla L(\theta)^\top \delta.
\]

\medskip

\noindent
\textbf{Choice of step size.}
\begin{itemize}
    \item Too large: instability
    \item Too small: slow convergence
\end{itemize}

\medskip

\noindent
\textbf{Variants.}
\begin{itemize}
    \item Momentum methods
    \item Adaptive methods (e.g., scaling by past gradients)
\end{itemize}

\medskip

\noindent
\textbf{Insight.}  
Gradient methods exploit local information to iteratively improve solutions.

\subsection{Stochastic Optimization}

In large-scale problems, computing the full gradient is often expensive.

\medskip

\noindent
\textbf{Stochastic gradient descent (SGD).}
\[
\theta_{t+1} = \theta_t - \eta \nabla \ell(f_\theta(x_i), y_i),
\]
where $(x_i, y_i)$ is a randomly selected data point.

\medskip

\noindent
\textbf{Mini-batch SGD.}
\[
\theta_{t+1} = \theta_t - \eta \frac{1}{|B|} \sum_{i \in B} \nabla \ell_i(\theta),
\]
where $B$ is a subset of the data.

\medskip

\noindent
\textbf{Properties.}
\begin{itemize}
    \item Lower computational cost per step
    \item Introduces noise into updates
\end{itemize}

\medskip

\noindent
\textbf{Effect of noise.}
\begin{itemize}
    \item Helps escape poor local minima
    \item Influences generalization behavior
\end{itemize}

\medskip

\noindent
\textbf{Tradeoff.}
\begin{itemize}
    \item Small batches: noisy but fast updates
    \item Large batches: stable but expensive
\end{itemize}

\medskip

\noindent
\textbf{Insight.}  
Stochasticity is not merely a computational trick; it plays a role in learning dynamics.

\subsection{Geometry of Loss Landscapes}

The objective function $L(\theta)$ defines a landscape over parameter space.

\medskip

\noindent
\textbf{Parameter space.}
\[
\theta \in \mathbb{R}^d
\]

\medskip

\noindent
\textbf{Loss landscape.}
\[
L : \mathbb{R}^d \to \mathbb{R}
\]

\medskip

\noindent
\textbf{Critical points.}
\begin{itemize}
    \item Local minima: $\nabla L(\theta) = 0$, positive curvature
    \item Saddle points: mixed curvature directions
\end{itemize}

\medskip

\noindent
\textbf{Convex vs non-convex.}
\begin{itemize}
    \item Convex problems: global minimum is unique
    \item Non-convex problems: many local minima and saddle points
\end{itemize}

\medskip

\noindent
\textbf{High-dimensional effects.}
\begin{itemize}
    \item Saddle points are more common than poor local minima
    \item Many directions may be nearly flat
\end{itemize}

\medskip

\noindent
\textbf{Flat vs sharp minima.}
\begin{itemize}
    \item Flat minima: stable under perturbations
    \item Sharp minima: sensitive to parameter changes
\end{itemize}

\medskip

\noindent
\textbf{Geometric intuition.}
\begin{itemize}
    \item Optimization follows paths along the landscape
    \item Curvature affects convergence speed
\end{itemize}

\medskip

\noindent
\textbf{Insight.}  
The geometry of the loss landscape determines both optimization behavior and generalization properties.

\subsection{A Unifying Perspective}

Learning as optimization brings together multiple aspects:

\begin{itemize}
    \item \textbf{Objective:} encodes data fit and regularization
    \item \textbf{Algorithm:} determines how solutions are found
    \item \textbf{Geometry:} shapes optimization trajectories
    \item \textbf{Stochasticity:} influences dynamics and generalization
\end{itemize}

\medskip

\noindent
\textbf{Core idea.}  
Learning is a dynamical process in parameter space driven by gradients.

\medskip

\noindent
\textbf{Key insight.}  
The outcome of learning depends not only on the objective, but also on the optimization algorithm used.

\subsection{Outlook}

In this lecture, we formulated learning as an optimization problem and studied:

\begin{itemize}
    \item objective functions for learning,
    \item gradient-based optimization,
    \item stochastic methods,
    \item geometric properties of loss landscapes.
\end{itemize}

\medskip

\noindent
These ideas lead to deeper questions:

\begin{itemize}
    \item Why do some models generalize better than others?
    \item How does optimization interact with model complexity?
    \item What role does the structure of the loss landscape play?
\end{itemize}

\medskip

\noindent
In the next lecture, we study generalization and overfitting, focusing on how models perform on unseen data.

\medskip

\noindent
\textbf{Guiding principle.}  
Learning is not only about finding a solution, but about finding a solution that generalizes.